\documentclass[10pt]{article}
\usepackage[T1]{fontenc}
\usepackage[utf8]{inputenc}
\usepackage[margin=0.7in]{geometry}
\usepackage{graphicx}
\usepackage{booktabs}
\usepackage{amsmath}
\usepackage{amssymb}
\usepackage{array}
\usepackage{caption}
\captionsetup{font=small}
\usepackage[hidelinks]{hyperref}
\usepackage{titlesec}
\titleformat*{\section}{\large\bfseries}
\titleformat*{\subsection}{\normalsize\bfseries}
\setlength{\parskip}{0.15em}
\graphicspath{{./}{figures/}}

\title{\textbf{Supplementary Information}}
\author{}
\date{}
\begin{document}
\maketitle
\sloppy
\section{Provenance of every reported value}
Values are tagged \emph{measured} (a real Boltz co-fold or REINVENT run), \emph{llm}, or \emph{heuristic}; the builder blocks any placeholder from a result.

\begin{table}[h]\centering\scriptsize\begin{tabular}{l l l}\toprule Value & Source & Detail \\ \midrule cytotoxin: derived rule + confidence & \textbf{llm} & 15 grounded exemplars, conf 0.784 \\
immunomodulator: derived rule + confidence & \textbf{llm} & 5 grounded exemplars, conf 0.821 \\
oligonucleotide: derived rule + confidence & \textbf{llm} & 3 grounded exemplars, conf 0.866 \\
cytotoxin: REINVENT-generated designs & \textbf{measured} & 25 assembled, best 0.6575 \\
immunomodulator: REINVENT-generated designs & \textbf{measured} & 25 assembled, best 0.7355 \\
oligonucleotide: REINVENT-generated designs & \textbf{measured} & 25 assembled, best 0.7374 \\
cytotoxin: held-out prediction & \textbf{llm} & verdict agree, conf 0.963 \\
immunomodulator: held-out prediction & \textbf{llm} & verdict agree, conf 0.828 \\
oligonucleotide: held-out prediction & \textbf{llm} & verdict disagree, conf 0.81 \\
Boltz cyto-Maleim-0: ipTM/Kd & \textbf{measured} & ipTM 0.745, Kd 8.9 nM \\
Boltz immu-Maleim-0: ipTM/Kd & \textbf{measured} & ipTM 0.859, Kd 46.48 nM \\
Boltz olig-DBCO-0: ipTM/Kd & \textbf{measured} & ipTM 0.883, Kd 108.69 nM \\
Boltz clinical-mcValCitPABC+MMAE: ipTM/Kd & \textbf{measured} & ipTM 0.791, Kd 35.23 nM \\
Synthesizability (Ertl SA\_Score) & \textbf{heuristic} & continuous SA\_Score, reverse-sigmoid to [ \\
Complexity index (step-count proxy) & \textbf{heuristic} & MW/chirality/ring/linkage formula; reporte \\
Mechanism stability score & \textbf{heuristic} & SMARTS liabilities: peptide/disulfide/male \\
Nearest clinical analogue (Tanimoto) & \textbf{heuristic} & Morgan fingerprint vs commercial linker se \\
ARC hypothesis co-fold: Kd & \textbf{measured} & Val-Cit ARC design, Kd 9.74 nM \\ \bottomrule\end{tabular}\end{table}
\section{Full reasoning chains and disagreement-aware confidence}
\subsection{Cytotoxin}
Retrieved 12 passages; grounded exemplars split 12 cleavable / 3 non-cleavable (unanimous); disagreement-aware confidence 0.60. Derived rule: cleavage=\textbf{protease-cleavable}, rigidity=semi-rigid. self\_confidence (LLM self-report, varied): 0.78. Rationale: The majority of approved and exemplified cytotoxin ADCs (MMAE, PBD dimer, PE38, Dxd) employ protease-cleavable dipeptide linkers based on the mc-ValCitPABC scaffold, which demonstrates high plasma stability and efficient lysosomal release. Variants incorporating acidic residues at the P3 position (GluValCit, AspValCit) further improve mouse plasma stability to 'very high', suggesting that charged extensions enhance serum stability without compromising clea
\subsection{Oligonucleotide}
Retrieved 12 passages; grounded exemplars split 0 cleavable / 3 non-cleavable (unanimous); disagreement-aware confidence 0.50. Derived rule: cleavage=\textbf{non-cleavable-rigid}, rigidity=rigid. self\_confidence (LLM self-report, varied): 0.55. Rationale: All three exemplars use non-cleavable thiol-maleimide linkers for siRNA conjugation, and among them the rigid linker sulfo-SMCC outperformed the flexible EMCS and the short AMAS in producing the desired DAR2 species. This indicates that for oligonucleotide payloads, a non-cleavable rigid linker with appropriate length is preferred to maximize conjugation efficiency and homogeneity. The sulfo group on SMCC likely aids aqueous solubility, which is important 
\subsection{Immunomodulator}
Retrieved 12 passages; grounded exemplars split 1 cleavable / 0 non-cleavable (unanimous); disagreement-aware confidence 0.83. Derived rule: cleavage=\textbf{protease-cleavable}, rigidity=flexible. self\_confidence (LLM self-report, varied): 0.25. Rationale: Only one exemplar provides concrete linker details: the imidazo-thienopyridine TLR7 agonist payload used a valine-citrulline dipeptide protease-cleavable linker with stochastic thiol-maleimide conjugation at interchain disulfide cysteines. The remaining four immunomodulator (TLR7/8) exemplars leave linker, conjugation, and cleavage unspecified, providing no corroborating or contradicting evidence. Therefore, the rule is extrapolated from a single well-char
\section{Held-out details}
\subsection{Cytotoxin}
Withheld: biomedicines-11-03080 (leaked: False). Predicted cleavage=\textbf{protease-cleavable}; verdict \textbf{AGREE}. Withheld paper states: \emph{Lysosomal cathepsin-B-cleavable Val-Cit/Val-Ala peptide linkers enable traceless cytotoxin release.} Blind rationale: Six of eight exemplars employ protease-cleavable linkers (dipeptide or tetrapeptide triggers with self-immolative PABC or AM spacers), establishing this as the dominant cleavage strategy for cytotoxic payloads. Plasma stability is a critical differentiator: the EVCit and cBuCit modifications over classical VCit demonstrate that extending or sterically shielding the peptide trigger markedly improves stability (half-li
\subsection{Immunomodulator}
Withheld: Immuno (leaked: False). Predicted cleavage=\textbf{protease-cleavable}; verdict \textbf{AGREE}. Withheld paper states: \emph{TLR7/8 ISAC linkers may be cleavable but demand high plasma stability to avoid systemic cytokine toxicity.} Blind rationale: Only one exemplar provides concrete linker details: the imidazo-thienopyridine TLR7 agonist ADC uses a Val-Cit dipeptide protease-cleavable linker with stochastic thiol-maleimide conjugation at interchain disulfide cysteines. This is consistent with the need for intracellular (endosomal/lysosomal) release of TLR7/8 agonists, which act on endosomal toll-like receptors, making cathepsin-cleavable dipeptide linkers a lo
\subsection{Oligonucleotide}
Withheld: siRNA, exploring-the-potentials-of-antibody-sirna-conjugates-in-tumor-cell-gene-silencing-without-cationic-assistance (leaked: False). Predicted cleavage=\textbf{protease-cleavable}; verdict \textbf{DISAGREE}. Withheld paper states: \emph{Antibody-oligonucleotide conjugates favour a rigid, NON-cleavable linker (sulfo-SMCC); a cleavable linker that sheds the polyanion is a liability.} Blind rationale: Among the exemplars, Val-Cit protease-cleavable linkers consistently enabled efficacious delivery of both siRNA and PMO payloads (DYNE-101, DYNE-251), while the Val-Arg variant failed, suggesting dipeptide identity matters and Val-Cit is preferred. Non-cleavable linkers (SMCC, maleimide) also showed activity (AOC-1001 achieving 80\% knockdown, MALAT1 ASO OTV), indicating oligonucleotide payloads can tolerate both clea
\section{All fifteen candidate dossiers}
\subsection{Cytotoxin}
\noindent\textbf{Cytotoxin \#1} \\ \texttt{CC(C)C(NC(=O)C(=O)NC(=O)CNC(=O)N1C(=O)C=CC1=O)C(=O)NC(CCCNC(N)=O)C(=O)} \\ This linker incorporates a Val-Cit dipeptide motif consistent with the protease-cleavable design rule preferred for MMAE payloads, enabling efficient cathepsin B-mediated lysosomal release. The maleimide conjugation handle and semi-rigid PABC-like self-immolative spacer (p-aminobenzyl alcohol) provi \\ Nearest clinical: mc-Val-Cit-PABC (Tc 0.654). Complexity idx 12, P(success) 0.14 (P(success) is depressed for peptidic linkers by the complexity-index proxy). \\ Route: Fmoc-SPPS assembly of Val-Cit dipeptide on Rink amide resin, with citrulline side-chain protection; Solution-phase coupling of p-aminobenzyl alcohol (PABOH) to Cit C-terminus via HATU-mediated amide bond formation; N-terminal glycine extension onto Val via Fmoc-Gly-OH coupling (H \\ Failure modes: Maleimide retro-Michael deconjugation in plasma leading to premature payload transfer to albumin (identified liability); Oxalamide bond hydrolysis under physiological conditions releasing the maleimid. Validation: Cathepsin B cleavage kinetics monitored by LC-MS/MS to confirm Val-Cit bond cleavage and complete 1,6-elimination of MMAE; Plasma stability assay (human, mouse, cynomolgus; 37\textdegree{}C, 7 days) with DAR trac. \\[0.4em]
\noindent\textbf{Cytotoxin \#2} \\ \texttt{CC(C)C(NC(=O)CNC(=O)C(=O)NC(=O)CN1C(=O)C=CC1=O)C(=O)NC(CCCNC(N)=O)C(=O} \\ This linker incorporates a Val-Cit dipeptide motif consistent with the protease-cleavable design rule preferred for MMAE payloads, enabling lysosomal cathepsin B-mediated release. The maleimide conjugation handle provides cysteine-selective attachment, and the semi-rigid scaffold with alternating am \\ Nearest clinical: mc-Val-Cit-PABC (Tc 0.711). Complexity idx 12, P(success) 0.14 (P(success) is depressed for peptidic linkers by the complexity-index proxy). \\ Route: Maleimide-glycine building block: alkylation of maleimide with bromoacetyl chloride followed by glycine amide coupling (HATU); Oxalamide formation: coupling of oxalyl chloride with glycine-maleimide intermediate under controlled stoichiometry; Dipeptide assembly: Fmoc-Val-OH coup \\ Failure modes: Maleimide retro-Michael deconjugation in plasma leading to off-target albumin transfer and reduced DAR over time; Oxalamide moiety may introduce hydrolytic instability under physiological pH, causing . Validation: Cathepsin B cleavage kinetics by LC-MS/MS monitoring MMAE release at pH 5.0 vs pH 7.4 over 24 h; Plasma stability assay in human, mouse, and cynomolgus monkey plasma measuring intact ADC DAR by HIC at. \\[0.4em]
\noindent\textbf{Cytotoxin \#3} \\ \texttt{CC(C)C(NC(=O)C(=O)NC(=O)CNC(=O)CN1C(=O)C=CC1=O)C(=O)NC(CCCNC(N)=O)C(=O} \\ This linker incorporates a Val-Cit dipeptide motif consistent with the protease-cleavable design rule preferred for MMAE payloads, enabling efficient cathepsin B-mediated lysosomal release. The maleimide conjugation handle and PABC-like self-immolative spacer (p-aminobenzyl alcohol) provide semi-rig \\ Nearest clinical: mc-Val-Cit-PABC (Tc 0.688). Complexity idx 12, P(success) 0.14 (P(success) is depressed for peptidic linkers by the complexity-index proxy). \\ Route: Maleimide-glycine conjugate formation: acylation of glycine with maleimidoacetic acid (EDC/HOBt); Elongation with second glycine unit via amide coupling (HATU/DIPEA); Coupling of Gly-Gly-maleimide fragment to Fmoc-Val-OH followed by Fmoc deprotection (standard Fmoc SPPS or soluti \\ Failure modes: Maleimide retro-Michael deconjugation in plasma leading to premature payload transfer to albumin (score 0.85 stability flag); Excessive linker length and flexibility from Gly-Gly spacer may reduce pla. Validation: Cathepsin B cleavage kinetics monitored by LC-MS/MS to confirm Val-Cit bond cleavage and MMAE release efficiency vs. mc-ValCitPABC benchmark; Plasma stability assay in human and mouse plasma (t1/2 mea. \\[0.4em]
\noindent\textbf{Cytotoxin \#4} \\ \texttt{CC(C)C(NC(=O)C(=O)NC(=O)CN(C)C(=O)CN1C(=O)C=CC1=O)C(=O)NC(CCCNC(N)=O)C} \\ This linker incorporates a Val-Cit dipeptide motif for protease-cleavable release of MMAE, consistent with the derived rule favoring cathepsin B-sensitive sequences for cytotoxin payloads. The maleimide conjugation handle and PABC-like self-immolative spacer (p-aminobenzyl alcohol) provide semi-rigi \\ Nearest clinical: mc-Val-Cit-PABC (Tc 0.654). Complexity idx 12, P(success) 0.14 (P(success) is depressed for peptidic linkers by the complexity-index proxy). \\ Route: Maleimide-glycine building block: N-alkylation of maleimide with bromoacetyl chloride followed by coupling to N-methylglycine (sarcosine) derivative; Oxalyl bridge formation: activation of oxalic acid mono-ester with EDC, coupling to valine-derived amine to install the oxalamide  \\ Failure modes: Maleimide retro-Michael deconjugation in plasma leading to premature payload transfer to serum albumin (score 0.85 stability flag); Oxalamide linkage susceptibility to non-specific hydrolysis under ph. Validation: Cathepsin B cleavage kinetics by LC-MS/MS monitoring MMAE release at pH 5.0 vs pH 7.4 over 24 h; Plasma stability assay in human, mouse, and cynomolgus monkey plasma measuring intact ADC by affinity c. \\[0.4em]
\noindent\textbf{Cytotoxin \#5} \\ \texttt{CC(C)C(NC(=O)CC(=O)NC(=O)C(=O)NC(=O)CN1C(=O)C=CC1=O)C(=O)NC(CCCNC(N)=O} \\ This linker incorporates a Val-Cit dipeptide motif consistent with the protease-cleavable design rule preferred for MMAE payloads, enabling efficient lysosomal cathepsin B-mediated release. The maleimide conjugation handle and semi-rigid amide backbone align with the mc-ValCitPABC clinical precedent \\ Nearest clinical: mc-Val-Cit-PABC (Tc 0.73). Complexity idx 12, P(success) 0.14 (P(success) is depressed for peptidic linkers by the complexity-index proxy). \\ Route: Maleimide-glycine building block: alkylation of glycine with maleic anhydride followed by ring closure to N-maleimidoacetyl-glycine; Oxalamide formation: coupling of N-maleimidoacetyl-glycine with oxalyl chloride mono-acid to install the oxalamide motif; Amide coupling of oxalami \\ Failure modes: Maleimide retro-Michael deconjugation in plasma leading to premature payload transfer to serum albumin; Oxalamide bond hydrolysis under physiological conditions causing linker fragmentation before tar. Validation: Cathepsin B cleavage kinetics by LC-MS/MS monitoring MMAE release at pH 5.0 vs pH 7.4 over 24 h; Human and mouse plasma stability assay (37\textdegree{}C, 96 h) with DAR tracking by HIC and intact-mass MS to quan. \\[0.4em]
\subsection{Immunomodulator}
\noindent\textbf{Immunomodulator \#1} \\ \texttt{CC(NC(=O)C(NC(=O)CC(N)C(=O)NCC(=O)N1C(=O)C=CC1=O)C(C)C)C(=O)Nc1ccc(CO)} \\ This linker employs a Val-Ala dipeptide protease-cleavable trigger consistent with the derived rule's cleavage preference for immunomodulator payloads, enabling cathepsin-mediated release in the endolysosome. The flexible alkyl spacer and PABC self-immolative unit match the 'flexible' rigidity prefe \\ Nearest clinical: mc-Val-Ala-PABC (Tc 0.58). Complexity idx 12, P(success) 0.14 (P(success) is depressed for peptidic linkers by the complexity-index proxy). \\ Route: Fmoc-SPPS assembly of Gly-Val-Ala tripeptide on Rink amide resin, cleaving to give H-Gly-Val-Ala-OH; Amide coupling of 4-aminobenzyl alcohol (PABC) with the tripeptide C-terminus using HATU/DIPEA; Activation of PABC benzylic alcohol as p-nitrophenyl carbonate, then coupling with  \\ Failure modes: Maleimide retro-Michael addition causing premature payload loss in plasma (thiol exchange with albumin Cys34); Steric hindrance from R848 imidazoquinoline ring system may impair PABC 1,6-elimination k. Validation: Cathepsin B cleavage assay (pH 5.0, 37\textdegree{}C) monitored by LC-MS/MS to confirm Val-Ala cleavage and free R848 release kinetics; Human and mouse plasma stability (37\textdegree{}C, 96 h) measuring intact ADC by HIC an. \\[0.4em]
\noindent\textbf{Immunomodulator \#2} \\ \texttt{CC(NC(=O)C(NC(=O)CS(=O)(=O)CC(=O)NCC(=O)N1C(=O)C=CC1=O)C(C)C)C(=O)Nc1c} \\ This linker employs a Val-Ala dipeptide protease-cleavable trigger consistent with the derived rule's cleavage preference for immunomodulator payloads, enabling selective intracellular release of R848 by cathepsin B. The flexible alkyl-sulfonyl spacer and PABC self-immolative unit match the 'flexibl \\ Nearest clinical: mc-Val-Ala-PABC (Tc 0.588). Complexity idx 12, P(success) 0.14 (P(success) is depressed for peptidic linkers by the complexity-index proxy). \\ Route: Sulfonylation: alkylation of thioglycolic acid with chloroacetamide-maleimide building block to install the sulfonyl-maleimide handle; Amide coupling of sulfonylacetyl acid with H-Ala-Val-OH dipeptide (HATU/DIPEA) to form the cleavable dipeptide segment; Synthesis of PABC-amine:  \\ Failure modes: Maleimide retro-Michael deconjugation in plasma leading to premature payload transfer to serum albumin; Sulfonyl linkage may introduce steric or electronic effects reducing cathepsin B cleavage effici. Validation: Cathepsin B cleavage kinetics monitored by LC-MS/MS to confirm Val-Ala cleavage and PABC self-immolation releasing free R848; Plasma stability assay (human, mouse) measuring intact ADC and free R848 a. \\[0.4em]
\noindent\textbf{Immunomodulator \#3} \\ \texttt{CC(NC(=O)C(NC(=O)C=CC(=O)NCC(=O)NCC(=O)N1C(=O)C=CC1=O)C(C)C)C(=O)Nc1cc} \\ This linker employs a Val-Ala dipeptide protease-cleavable trigger consistent with the derived rule's cleavage preference for immunomodulator payloads, enabling cathepsin-mediated release in the endolysosome. The flexible PEG-like glycine spacers and linear amide backbone align with the 'flexible' r \\ Nearest clinical: mc-Val-Ala-PABC (Tc 0.549). Complexity idx 12, P(success) 0.14 (P(success) is depressed for peptidic linkers by the complexity-index proxy). \\ Route: Maleimide-glycine building block synthesis: N-maleimidoacetic acid preparation via maleic anhydride ring closure with glycine; Amide coupling of maleimide-Gly-OH with H-Gly-OH spacer (HATU/DIPEA); Maleic acid monoamide formation: conjugation of maleimide-Gly-Gly with fumaric acid \\ Failure modes: Maleimide retro-Michael deconjugation in plasma leading to premature payload transfer to serum albumin and off-target immune activation; 12-step linear synthesis reduces overall yield and scalability,. Validation: Cathepsin B cleavage kinetics by LC-MS/MS monitoring R848 release at pH 5.0 vs pH 7.4 over 24 h; Human and mouse plasma stability assay (37\textdegree{}C, 0--96 h) with DAR tracking by HIC and free R848 quantifica. \\[0.4em]
\noindent\textbf{Immunomodulator \#4} \\ \texttt{CC(NC(=O)C(NC(=O)CCC(=O)NCC(=O)NCC(=O)N1C(=O)C=CC1=O)C(C)C)C(=O)Nc1ccc} \\ This linker employs a Val-Ala dipeptide protease-cleavable trigger consistent with the derived rule's cleavage preference for immunomodulator payloads, enabling cathepsin-mediated release in the endolysosomal compartment. The flexible Glu-Gly-Gly spacer between the maleimide conjugation handle and t \\ Nearest clinical: mc-Val-Ala-PABC (Tc 0.646). Complexity idx 12, P(success) 0.14 (P(success) is depressed for peptidic linkers by the complexity-index proxy). \\ Route: Fmoc solid-phase peptide synthesis of Glu(OtBu)-Gly-Gly-Val-Ala pentapeptide fragment with N-terminal maleimidocaproyl capping; Solution-phase synthesis of 4-aminobenzyl alcohol (PABC) carbonate or carbamate linkage to R848 exocyclic amine; Amide coupling of C-terminal Ala of dip \\ Failure modes: Maleimide retro-Michael deconjugation in plasma leading to premature payload transfer to serum albumin and off-target TLR activation; Incomplete PABC 1,6-elimination under lysosomal conditions resulti. Validation: Cathepsin B cleavage kinetics monitored by LC-MS/MS to confirm Val-Ala cleavage and quantify free R848 release efficiency; Plasma stability assay in human and mouse plasma (37\textdegree{}C, 96 h) measuring intac. \\[0.4em]
\noindent\textbf{Immunomodulator \#5} \\ \texttt{CC(NC(=O)C(NC(=O)C(N)CCC(=O)NCC(=O)N1C(=O)C=CC1=O)C(C)C)C(=O)Nc1ccc(CO} \\ This linker employs a Val-Ala dipeptide protease-cleavable trigger consistent with the derived rule's cleavage preference for immunomodulator payloads, enabling cathepsin-mediated release in the lysosome. The flexible alkyl spacer and PABC self-immolative unit match the 'flexible' rigidity preferenc \\ Nearest clinical: mc-Val-Ala-PABC (Tc 0.563). Complexity idx 12, P(success) 0.14 (P(success) is depressed for peptidic linkers by the complexity-index proxy). \\ Route: Fmoc-SPPS assembly of Glu($\gamma$-glycine)-Val-Ala tripeptide on resin, cleaving to give H-Glu($\gamma$-Gly)-Val-Ala-OH; Amide coupling of tripeptide C-terminus with 4-aminobenzyl alcohol (PABC) using HATU/DIPEA; Activation of PABC benzyl alcohol as p-nitrophenyl carbonate, then displacement  \\ Failure modes: Maleimide retro-Michael deconjugation in plasma leading to premature payload transfer to serum albumin; Incomplete or slow cathepsin B/L cleavage of Val-Ala dipeptide due to steric effects of the bran. Validation: Cathepsin B cleavage kinetics monitored by LC-MS/MS to confirm Val-Ala dipeptide processing and free R848 release; Plasma stability assay (human, mouse) measuring intact ADC and free payload over 96 h. \\[0.4em]
\subsection{Oligonucleotide}
\noindent\textbf{Oligonucleotide \#1} \\ \texttt{NCC1CCC(NC(=O)N2CCN(c3ccc(S(=O)(=O)N4Cc5ccccc5C\#Cc5ccccc54)cn3)CC2)CC1} \\ This linker employs a non-cleavable urea-piperazine backbone with a fused dibenzocyclooctyne (DBCO)-reactive sulfonamide motif, directly satisfying the derived rule's preference for non-cleavable-rigid architectures for siRNA payloads. The rigid aromatic alkyne scaffold and constrained piperazine ri \\ Nearest clinical: MCC (non-cleavable) (Tc 0.145). Complexity idx 6, P(success) 0.65 (P(success) is depressed for peptidic linkers by the complexity-index proxy). \\ Route: Sulfonamide formation: react 2-aminodibenzazocine (iminodibenzyl alkyne) with pyridine-4-sulfonyl chloride bearing a piperazine handle; SNAr or Buchwald coupling: install piperazine onto chloropyridinesulfonamide scaffold; Urea formation: react piperazine nitrogen with 4-aminocyc \\ Failure modes: Steric congestion around the sulfonamide-dibenzocyclooctyne junction may impair SPAAC click conjugation kinetics with azide-modified siRNA; The extended aromatic system may increase hydrophobicity bey. Validation: DAR analysis by HIC and native SEC to assess conjugation homogeneity and confirm DAR2 enrichment; Plasma stability assay (human and mouse, 37\textdegree{}C, 7 days) with LC-MS/MS quantification of intact conjugat. \\[0.4em]
\noindent\textbf{Oligonucleotide \#2} \\ \texttt{NCC1CCC(NC(=O)c2ccc(NC(=O)N3Cc4ccccc4C\#Cc4ccccc43)nc2)CC1} \\ This linker satisfies the non-cleavable-rigid design rule for oligonucleotide payloads by incorporating a urea-bridged diarylacetylene core that enforces rigidity analogous to the sulfo-SMCC cyclohexane scaffold, promoting homogeneous DAR2 conjugation. The DBCO handle enables strain-promoted azide-a \\ Nearest clinical: MCC (non-cleavable) (Tc 0.158). Complexity idx 6, P(success) 0.65 (P(success) is depressed for peptidic linkers by the complexity-index proxy). \\ Route: Sonogashira coupling of 2-iodobenzylamine derivative with 2-ethynylaniline to form the diarylacetylene DBCO-like core; Intramolecular cyclization to close the dibenzoazacyclooctyne (DBCO) ring system via ring-closing amination; Urea formation: reaction of the DBCO-amine with 5-am \\ Failure modes: DBCO cyclooctyne strain may lead to non-specific thiol reactivity with serum albumin, reducing effective siRNA delivery; Extended rigid aromatic core increases hydrophobicity, potentially causing aggr. Validation: DAR analysis by HIC and native SEC-MALS to confirm homogeneous DAR2 conjugation with siRNA; Plasma stability assay (human, mouse, 37\textdegree{}C, 72h) with LC-MS/MS quantification of intact conjugate vs. free s. \\[0.4em]
\noindent\textbf{Oligonucleotide \#3} \\ \texttt{NCC1CCC(Nc2ccc(NC(=O)N3CCN(C(=O)N4Cc5ccccc5C\#Cc5ccccc54)CC3)cn2)CC1} \\ This linker employs a non-cleavable urea-piperazine-urea backbone with a fused isoindoline-alkyne rigid scaffold, directly matching the derived rule's preference for non-cleavable-rigid architectures for siRNA payloads. The DBCO moiety provides strain-promoted azide-alkyne cycloaddition (SPAAC) conj \\ Nearest clinical: MCC (non-cleavable) (Tc 0.152). Complexity idx 8, P(success) 0.4 (P(success) is depressed for peptidic linkers by the complexity-index proxy). \\ Route: Buchwald-Hartley amination of 4-amino-2-chloropyridine with trans-4-aminocyclohexylamine to install the diaminopyridine-cyclohexane fragment; Urea formation between the pyridine amine and piperazine mono-Boc using CDI or triphosgene; Boc deprotection of piperazine nitrogen with T \\ Failure modes: 8-step linear synthesis reduces overall yield and scalability, increasing cost-of-goods for clinical development; Low Tanimoto similarity (0.152) to validated MCC linker suggests uncertain translation. Validation: DAR analysis by HIC and native SEC-MALS to assess conjugation homogeneity and confirm DAR2 preference consistent with sulfo-SMCC benchmark; Plasma stability assay (human, mouse, 37\textdegree{}C, 7 days) with LC-. \\[0.4em]
\noindent\textbf{Oligonucleotide \#4} \\ \texttt{NCC1CCC(NC(=O)c2cnc(-c3cnc(N4Cc5ccccc5C\#Cc5ccccc54)nc3)nc2)CC1} \\ This linker satisfies the non-cleavable-rigid design rule for siRNA payloads by incorporating a fused bipyrimidine-dihydroisoindole core that enforces conformational rigidity, minimizing entropic penalties during conjugation and promoting homogeneous DAR species. The DBCO handle enables strain-promo \\ Nearest clinical: MCC (non-cleavable) (Tc 0.16). Complexity idx 5, P(success) 0.85 (P(success) is depressed for peptidic linkers by the complexity-index proxy). \\ Route: Suzuki coupling of 5-bromopyrimidine-2-carboxylic acid with 5-(4,4,5,5-tetramethyl-1,3,2-dioxaborolan-2-yl)pyrimidine-2-amine to form the bipyrimidine core; Buchwald-Hartley N-arylation of the 2-aminopyrimidine with 2-ethynyl-iodobenzene to install the alkyne-aniline moiety; Intr \\ Failure modes: Poor aqueous solubility of the extended aromatic scaffold may limit conjugation efficiency with hydrophilic siRNA in aqueous buffers, unlike the sulfo-SMCC benchmark which carries a solubilizing sulfo. Validation: SPAAC conjugation kinetics with azide-modified siRNA monitored by HPLC and MALDI-TOF to confirm click efficiency and DAR distribution (target DAR2); Plasma stability assay in human and mouse plasma at. \\[0.4em]
\noindent\textbf{Oligonucleotide \#5} \\ \texttt{NCC1CCC(NS(=O)(=O)N2CCN(C(=O)c3ccc(N4Cc5ccccc5C\#Cc5ccccc54)nc3)CC2)CC1} \\ This linker satisfies the non-cleavable-rigid design rule for oligonucleotide payloads by incorporating a DBCO click-chemistry handle for strain-promoted azide-alkyne cycloaddation (SPAAC) conjugation, ensuring stable triazole linkage that is non-cleavable in circulation. Rigidity is enforced throug \\ Nearest clinical: MCC (non-cleavable) (Tc 0.163). Complexity idx 5, P(success) 0.85 (P(success) is depressed for peptidic linkers by the complexity-index proxy). \\ Route: Synthesis of 4-aminocyclohexyl amine core via reductive amination or Boc-protection/deprotection of trans-1,4-diaminocyclohexane; Sulfonylation of cyclohexylamine nitrogen with piperazine-1-sulfonyl chloride to install the sulfonamide-piperazine unit; Amide coupling of piperazine \\ Failure modes: DBCO moiety may undergo premature reaction with biological thiols or nucleophiles in formulation buffer, reducing conjugation efficiency; The extended aromatic DBCO system may increase lipophilicity b. Validation: DAR analysis by HIC and native SEC-MALS after SPAAC conjugation with azide-functionalized siRNA to assess conjugation efficiency and homogeneity; Plasma stability assay (human, mouse, 37\textdegree{}C, 7 days) wi. \\[0.4em]
\section{Boltz-2 co-folds (linker + payload co-present, cathepsin B)}
\begin{table}[h]\centering\small\begin{tabular}{l c c}\toprule Co-fold (co-present) & ipTM & Pred.\ $K_d$ (nM) \\ \midrule cyto-Maleim-0 & 0.745 & 8.9 \\
immu-Maleim-0 & 0.859 & 46.48 \\
olig-DBCO-0 & 0.883 & 108.69 \\
clinical-mcValCitPABC+MMAE & 0.791 & 35.23 \\ \bottomrule\end{tabular}\end{table}
\section{Literature corpus (30 documents the agent read)}
The retrieval and reasoning operated over the following primary-literature documents. Identifiers/titles are the sources as ingested (filenames, DOIs, or journal handles); they are the ground truth for every grounded exemplar cited above.

\begin{enumerate}\setlength{\itemsep}{0pt}\small
\item 1-s2.0-S0960894X23002263-main (Elsevier ScienceDirect)
\item 1-s2.0-S104084282500246X-main (Elsevier ScienceDirect)
\item 1-s2.0-S104366182400104X-main (Elsevier ScienceDirect)
\item 10555 2024 Article 10231 (Springer)
\item 1441
\item Advanced Science - 2025 - Long - Novel Quaternary Ammonium SaltLinked STING Agonist AntibodyDrug Conjugate  Synergistic
\item ChemMedChem - 2025 - Lei - Linker Design for the Antibody Drug Conjugates  A Comprehensive Review
\item ChemSocRev (Chem.\textbackslash\{\} Soc.\textbackslash\{\} Rev.)
\item ISAC
\item Immuno
\item PIIS0923753426001493 (Elsevier/Cell Press)
\item PIIS2162253125000848 (Elsevier/Cell Press)
\item biomedicines-11-03080 (Biomedicines (MDPI))
\item cir-24-0103 (Cancer Immunol.\textbackslash\{\} Res.)
\item d2cs00446a (Chem.\textbackslash\{\} Soc.\textbackslash\{\} Rev., DOI 10.1039/d2cs00446a)
\item Discovery and optimization of a sting agonist platform for application in antibody drug conjugates
\item Dissecting the efficacy and immunogenicity of tlr7 agonist antibody conjugates through the lens of fc effector function
\item Exploring the potentials of antibody sirna conjugates in tumor cell gene silencing without cationic assistance
\item fphar-14-1320524 (Front.\textbackslash\{\} Pharmacol.)
\item Identification of a novel linker enabling the bioconjugation of a cyclic dinucleotide for the sting antibody drug
\item jitc-13-8 (J.\textbackslash\{\} Immunother.\textbackslash\{\} Cancer)
\item Metabolic stability and targeted delivery of oligonucleotides advancing rna therapeutics beyond the liver
\item pharmacology
\item S41434 026 00621 5
\item siRNA
\item songli2021
\item Structure activity relationship of antibody oligonucleotide 3 conjugates evaluating bioconjugation strategies for
\item Structure activity relationship of antibody oligonucleotide conjugates evaluating bioconjugation strategies for
\item tbag022
\item Thiol specific linkers for the synthesis of oligonucleotide conjugates via metal free thiol ene click reaction
\end{enumerate}
\end{document}